% concatenated from main-4.tex
\documentclass[12pt]{amsart}

\usepackage[all]{xy}
\usepackage{fullpage}
\usepackage{latexsym}
\usepackage{amsmath}
\usepackage{amsfonts}
\usepackage{amssymb}
\usepackage{amsthm}
\usepackage{eucal}
\usepackage{enumerate,yfonts}
\usepackage{mathrsfs}
\usepackage{graphicx}
\usepackage{graphics}
\usepackage{epstopdf}
\usepackage{amscd}
\usepackage{bbm}
\usepackage{hyperref}
\usepackage{url}
\usepackage{color}
\usepackage{bbm}
\usepackage{cancel}
\usepackage{enumerate}
\usepackage{amsmath,amsthm}
\usepackage{amssymb}
\usepackage{epsfig}
\usepackage{pstricks}
\usepackage{xy}
\usepackage{xypic}
\usepackage{epigraph}

%%%






\newtheorem{thm}{Theorem}[section]
\newtheorem{corollary}[thm]{Corollary}
\newtheorem{lemma}[thm]{Lemma}
\newtheorem{lem}[thm]{Lemma}
\newtheorem{proposition}[thm]{Proposition}
\newtheorem{prop}[thm]{Proposition}
\newtheorem{conjecture}[thm]{Conjecture}
\newtheorem{thm-dfn}[thm]{Theorem-Definition}
\newtheorem{claim}{Claim}[section]


\newtheorem*{coro}{Corollary} %%%% for unnumbered statements

\theoremstyle{definition}
\newtheorem{definition}[thm]{Definition}

\numberwithin{equation}{section}

\theoremstyle{remark}
\newtheorem{remark}{Remark}[section]
\newtheorem{example}[remark]{Example}

\newtheorem{rmk}[remark]{Remark}
\newtheorem{exam}[remark]{Example}



\newcommand{\cC}{{\mathcal C}}% one could define new command for convenience
\newcommand{\fg}{{\mathfrak g}}
\newcommand{\ft}{{\mathfrak t}}
\newcommand{\fl}{{\mathfrak l}}
\newcommand{\fb}{{\mathfrak b}}
\newcommand{\fu}{{\mathfrak u}}
\newcommand{\fp}{{\mathfrak p}}
\newcommand{\fz}{{\mathfrak z}}
\newcommand{\ff}{{\mathfrak F}}
\newcommand{\fa}{{\mathfrak a}}
\newcommand{\fC}{{\mathfrak C}}
\newcommand{\fX}{{\mathfrak X}}
\newcommand{\fc}{{\mathfrak c}}
\newcommand{\fk}{{\mathfrak k}}





\newcommand{\rZ}{{\mathrm Z}}
\newcommand{\rW}{{\mathrm W}}
\newcommand{\rU}{{\mathrm U}}
\newcommand{\rUg}{\mathrm U(\fg)}
\newcommand{\rUl}{\mathrm U(\fl)}

\newcommand{\hL}{\widehat{\lambda+\Lambda}}
\newcommand{\ol}{\overline{\lambda}}
\newcommand{\hl}{\hat{\lambda}}
\newcommand{\hol}{\widehat{\lambda+\Lambda}}

\newcommand{\bC}{{\mathbb C}}
\newcommand{\bX}{{\mathbb X}}
\newcommand{\bG}{{\mathbb G}}
\newcommand{\bZ}{{\mathbb Z}}
\newcommand{\bB}{{\mathbb B}}
\newcommand{\bT}{{\mathbb T}}
\newcommand{\bH}{{\mathbb H}}


\newcommand{\tD}{\widetilde{\D}}

\newcommand{\mD}{\mathcal{D}}
\newcommand{\mS}{\mathcal{S}}
\newcommand{\mI}{\mathcal{I}}
\newcommand{\mY}{\mathcal{Y}}
\newcommand{\mX}{\mathcal{X}}
\newcommand{\mE}{\mathcal{E}}
\newcommand{\mF}{\mathcal{F}}
\newcommand{\mQ}{\mathcal{Q}}
\newcommand{\mJ}{J}
\newcommand{\mA}{\mathcal{A}}
\newcommand{\mM}{\mathcal{M}}
\newcommand{\mT}{\mathcal{T}}
\newcommand{\mO}{\mathcal{O}}
\newcommand{\mL}{\mathcal{L}}
\newcommand{\mH}{\mathcal{H}}
\newcommand{\mP}{\mathcal{P}}
\newcommand{\mad }{\mathcal{A}^{\diamond}}
\newcommand{\mG}{\mathcal{G}}
\newcommand{\mN}{\mathcal{N}}
\newcommand{\mB}{\mathcal{B}}

\newcommand{\sX}{\mathscr{X}}
\newcommand{\sY}{\mathscr{Y}}
\newcommand{\sS}{\mathscr{S}}
\newcommand{\sG}{\mathscr{G}}
\newcommand{\sB}{\mathscr{B}}
\newcommand{\sP}{\mathscr{P}}
\newcommand{\sT}{\mathscr{T}}
\newcommand{\sA}{\mathscr{A}}
\newcommand{\sM}{\mathscr{M}}
\newcommand{\scP}{\mathscr{P}}
\newcommand{\sH}{\mathscr{H}}
\newcommand{\sL}{\mathscr{L}}
\newcommand{\sQ}{\mathscr{Q}}
\newcommand{\sC}{\mathscr{C}}
\newcommand{\sD}{\mathscr{D}}

\newcommand{\on}{\operatorname}
\newcommand{\D}{\on{D}}
\newcommand{\h}{\on{H}}
\newcommand{\tl}{t^{\lambda}}
\newcommand{\tU}{\widetilde U}
\newcommand{\tX}{\widetilde X}

\newcommand{\oP}{\overline{P}}
\newcommand{\oN}{\overline{N}}
\newcommand{\uP}{\underline{P}}

\newcommand{\ra}{\rightarrow}
\newcommand{\la}{\leftarrow}
\newcommand{\gl}{\geqslant}
\newcommand{\s}{\textbackslash}
\newcommand{\bs}{\backslash}
\newcommand{\is}{\simeq}

\newcommand{\Higgs}{{\on{Higgs}}}
\newcommand{\Loc}{\on{Loc}}
\newcommand{\Bun}{\on{Bun}}
\newcommand{\Sect}{\on{Sect}}
\newcommand{\tC}{\widetilde C}

\newcommand{\quash}[1]{}  %%Anything in \quash is ignored
\newcommand{\nc}{\newcommand}



\newcommand{\fraka}{{\mathfrak a}}
\newcommand{\frakb}{{\mathfrak b}}
\newcommand{\frakc}{{\mathfrak c}}
\newcommand{\frakd}{{\mathfrak d}}
\newcommand{\frake}{{\mathfrak e}}
\newcommand{\frakf}{{\mathfrak f}}
\newcommand{\frakg}{{\mathfrak g}}
\newcommand{\frakh}{{\mathfrak h}}
\newcommand{\fraki}{{\mathfrak i}}
\newcommand{\frakj}{{\mathfrak j}}
\newcommand{\frakk}{{\mathfrak k}}
\newcommand{\frakl}{{\mathfrak l}}
\newcommand{\frakm}{{\mathfrak m}}
\newcommand{\frakn}{{\mathfrak n}}
\newcommand{\frako}{{\mathfrak o}}
\newcommand{\frakp}{{\mathfrak p}}
\newcommand{\frakq}{{\mathfrak q}}
\newcommand{\frakr}{{\mathfrak r}}
\newcommand{\fraks}{{\mathfrak s}}
\newcommand{\frakt}{{\mathfrak t}}
\newcommand{\fraku}{{\mathfrak u}}
\newcommand{\frakv}{{\mathfrak v}}
\newcommand{\frakw}{{\mathfrak w}}
\newcommand{\frakx}{{\mathfrak x}}
\newcommand{\fraky}{{\mathfrak y}}
\newcommand{\frakz}{{\mathfrak z}}
\newcommand{\frakA}{{\mathfrak A}}
\newcommand{\frakB}{{\mathfrak B}}
\newcommand{\frakC}{{\mathfrak C}}
\newcommand{\frakD}{{\mathfrak D}}
\newcommand{\frakE}{{\mathfrak E}}
\newcommand{\frakF}{{\mathfrak F}}
\newcommand{\frakG}{{\mathfrak G}}
\newcommand{\frakH}{{\mathfrak H}}
\newcommand{\frakI}{{\mathfrak I}}
\newcommand{\frakJ}{{\mathfrak J}}
\newcommand{\frakK}{{\mathfrak K}}
\newcommand{\frakL}{{\mathfrak L}}
\newcommand{\frakM}{{\mathfrak M}}
\newcommand{\frakN}{{\mathfrak N}}
\newcommand{\frakO}{{\mathfrak O}}
\newcommand{\frakP}{{\mathfrak P}}
\newcommand{\frakQ}{{\mathfrak Q}}
\newcommand{\frakR}{{\mathfrak R}}
\newcommand{\frakS}{{\mathfrak S}}
\newcommand{\frakT}{{\mathfrak T}}
\newcommand{\frakU}{{\mathfrak U}}
\newcommand{\frakV}{{\mathfrak V}}
\newcommand{\frakW}{{\mathfrak W}}
\newcommand{\frakX}{{\mathfrak X}}
\newcommand{\frakY}{{\mathfrak Y}}
\newcommand{\frakZ}{{\mathfrak Z}}
\newcommand{\bbA}{{\mathbb A}}
\newcommand{\bbB}{{\mathbb B}}
\newcommand{\bbC}{{\mathbb C}}
\newcommand{\bbD}{{\mathbb D}}
\newcommand{\bbE}{{\mathbb E}}
\newcommand{\bbF}{{\mathbb F}}
\newcommand{\bbG}{{\mathbb G}}
\newcommand{\bbH}{{\mathbb H}}
\newcommand{\bbI}{{\mathbb I}}
\newcommand{\bbJ}{{\mathbb J}}
\newcommand{\bbK}{{\mathbb K}}
\newcommand{\bbL}{{\mathbb L}}
\newcommand{\bbM}{{\mathbb M}}
\newcommand{\bbN}{{\mathbb N}}
\newcommand{\bbO}{{\mathbb O}}
\newcommand{\bbP}{{\mathbb P}}
\newcommand{\bbQ}{{\mathbb Q}}
\newcommand{\bbR}{{\mathbb R}}
\newcommand{\bbS}{{\mathbb S}}
\newcommand{\bbT}{{\mathbb T}}
\newcommand{\bbU}{{\mathbb U}}
\newcommand{\bbV}{{\mathbb V}}
\newcommand{\bbW}{{\mathbb W}}
\newcommand{\bbX}{{\mathbb X}}
\newcommand{\bbY}{{\mathbb Y}}
\newcommand{\bbZ}{{\mathbb Z}}
\newcommand{\calA}{{\mathcal A}}
\newcommand{\calB}{{\mathcal B}}
\newcommand{\calC}{{\mathcal C}}
\newcommand{\calD}{{\mathcal D}}
\newcommand{\calE}{{\mathcal E}}
\newcommand{\calF}{{\mathcal F}}
\newcommand{\calG}{{\mathcal G}}
\newcommand{\calH}{{\mathcal H}}
\newcommand{\calI}{{\mathcal I}}
\newcommand{\calJ}{{\mathcal J}}
\newcommand{\calK}{{\mathcal K}}
\newcommand{\calL}{{\mathcal L}}
\newcommand{\calM}{{\mathcal M}}
\newcommand{\calN}{{\mathcal N}}
\newcommand{\calO}{{\mathcal O}}
\newcommand{\calP}{{\mathcal P}}
\newcommand{\calQ}{{\mathcal Q}}
\newcommand{\calR}{{\mathcal R}}
\newcommand{\calS}{{\mathcal S}}
\newcommand{\calT}{{\mathcal T}}
\newcommand{\calU}{{\mathcal U}}
\newcommand{\calV}{{\mathcal V}}
\newcommand{\calW}{{\mathcal W}}
\newcommand{\calX}{{\mathcal X}}
\newcommand{\calY}{{\mathcal Y}}
\newcommand{\calZ}{{\mathcal Z}}


\nc{\al}{{\alpha}} \nc{\be}{{\beta}} \nc{\ga}{{\gamma}}
\nc{\ve}{{\varepsilon}} \nc{\Ga}{{\Gamma}} %\nc{\la}{{\lambda}}
\nc{\La}{{\Lambda}}

\newcommand{\Dmod}{\on{D-mod}}

\nc{\ad }{{\on{ad }}}

\nc{\aff}{{\on{aff}}} \nc{\Aff}{{\mathbf{Aff}}}
\newcommand{\Aut}{{\on{Aut}}}
\newcommand{\Out}{{\on{Out}}}
%\nc{\Bun}{{\on{Bun}}}
\newcommand{\cha}{{\on{char}}}
\nc{\der}{{\on{der}}}
\newcommand{\Der}{{\on{Der}}}
\nc{\diag}{{\on{diag}}}
\newcommand{\End}{{\on{End}}}
\nc{\Fl}{{\calF\ell}}
\newcommand{\Gal}{{\on{Gal}}}
\newcommand{\Gr}{{\on{Gr}}}
\nc{\Hg}{{\on{Higgs}}}
\newcommand{\Hom}{{\on{Hom}}}
\newcommand{\id}{{\on{id}}}
\nc{\Id}{{\on{Id}}}
\newcommand{\ind}{{\on{ind}}}
\nc{\Ind}{{\on{Ind}}}
\newcommand{\Lie}{{\on{Lie}}}
\nc{\Op}{{\on{Op}}}
\newcommand{\Pic}{{\on{Pic}}}
\newcommand{\pr}{{\on{pr}}}
\newcommand{\Res}{{\on{Res}}}
\nc{\res}{{\on{res}}}
%\newcommand{\s}{{\on{sc}}}
\newcommand{\sgn}{{\on{sgn}}}
\newcommand{\Spec}{{\on{Spec}}}
\newcommand{\St}{{\on{St}}}
\nc{\tr}{{\on{tr}}}
\newcommand{\Tr}{{\on{Tr}}}
\newcommand{\Mod}{{\mathrm{-Mod}}}
\newcommand{\GL}{{\on{GL}}}
%\newcommand{\gl}{{\frakg\frakl}}
\nc{\GSp}{{\on{GSp}}} \nc{\GU}{{\on{GU}}} \nc{\SL}{{\on{SL}}}
\nc{\SU}{{\on{SU}}} \nc{\SO}{{\on{SO}}}

\nc{\nh}{{\Loc_{J^p}(\tau')}}
\nc{\bnh}{{\Loc_{\breve J^p}(\tau')}}
\newcommand{\bGr}{{\overline{\Gr}}}
\nc{\bU}{{\overline{U}}} \nc{\IC}{{\on{IC}}}
\newcommand{\Nm}{{\on{Nm}}}
\newcommand{\ppart}{(\!(t)\!)}
\newcommand{\pparu}{(\!(u)\!)}

\def\xcoch{\mathbb{X}_\bullet}
\def\xch{\mathbb{X}^\bullet}

\newcommand{\Chow}{{\on{Chow}}}
\newcommand{\AJ}{{\on{AJ}}}

\newcommand{\beqn}{\begin{equation*}}
\newcommand{\eeqn}{\end{equation*}}

\newcommand{\beq}{\begin{equation}}
\newcommand{\eeq}{\end{equation}}

\newcommand{\wsM}{\widetilde{\mathscr M}}
\newcommand{\CC}[1]{{\color{red!40!black}[CC: #1]}}
\newcommand{\Irr}{\operatorname{Irr}}
\newcommand{\Rep}{\operatorname{Rep}}
\newcommand{\Grs}{G^{\mathrm{rs}}}




\quash{

\setlength{\parskip}{2ex}
\setlength{\oddsidemargin}{0in}
\setlength{\evensidemargin}{0in}
\setlength{\textwidth}{6.5in}
\setlength{\topmargin}{-0.15in}
\setlength{\textheight}{8.6in}


\topmargin-0.5cm \textheight22cm \oddsidemargin1.2cm \textwidth14cm}
\begin{document}
\title{Admissibility of Bernstein centers}

\author{Tsao-Hsien Chen}
        \address{School of Mathematics, University of Minnesota, Twin cities, Minneapolis, MN 55455 }
         \email{chenth@umn.edu}
         \author{Cheng-Chiang Tsai}
        \address{Institute of Mathematics, Academia Sinica, 6F, Astronomy-Mathematics Building,
No. 1, Sec. 4, Roosevelt Road, Taipei, Taiwan\vskip.2cm
also National Sun Yat-Sen University and National Taiwan University}
         \email{chchtsai@as.edu.tw}
\thanks{}
\thanks{}

\maketitle     
\begin{abstract}
We provide a criterion for determining when elements of the Bernstein center of a totally disconnected locally compact group are admissible invariant distributions in the sense of Harish-Chandra \cite{HC99}. As a consequence, we deduce the local integrability results for elements of bounded depth or Bernstein supports in the Bernstein centers of reductive $p$-adic groups. Our methods apply uniformly to both complex and mod-$\ell$ coefficients, generalizing results of Moy and Tadi\'{c} \cite{MT02} in the complex case.



\end{abstract}
\section{Bernstein centers}
We recall some basic facts about the Bernstein center, following \cite{BD84}. Let $G$ be a totally disconnected locally compact group. Fix a prime $p$, and assume that $G$ contains a compact open pro-$p$ subgroup. Let $k$ be an algebraically closed field of characteristic $\ell\neq p$. Denote by $(H(G),*)$ the $k$-algebra of locally constant $k$-valued measures on $G$ with compact support. Let $R(G)$ be the category of smooth $k$-representations of $G$. Let $\on{Irr}(G)$ be the set of irreducible $k$-representations of $G$.


 There is an equivalence between 
$R(G)$ and the category of non-degenerate $H(G)$-modules, that is,
 $H(G)$-modules $V$ satisfying $H(G)\cdot V=V$.
The action of $f\in H(G)$ on a representation 
$(\pi,V_\pi)\in R(G)$
is given by
\[\pi(f): V_\pi\to V_\pi,\ \pi(f)(v)=\int_G f(g)\pi(g)(v).\]
Let \[Z(G)=\on{End}_{R(G)}(\on{Id})\] denote the Bernstein center of $G$.
Let $D(G)$ be the space of $k$-valued distributions on $G$, and let $D(G)^G$ denote the subspace of invariant distributions.
There is a canonical embedding \[\iota:Z(G)\to D(G)^G\]
given by $\iota(z)(h)=z(h^*)(\operatorname{id})$ where $h^*(g)=h(g^{-1})$ and $z(h^*)$ is the action of $z$ on $H(G)$ viewed as an $H(G)$-module.
The image of $\iota$ is the space of essentially compact invariant distributions 
\[D(G)^G_{ec}=\{f\in D(G)^G\;|\;f*h\in H(G)\text{\ \ for any\ \ } h\in H(G) \}.\]
For any $z\in Z(G)$,  we write
\[f_z=\iota(z)\in D(G)^G_{ec}\] the corresponding essentially compact invariant distribution.

Given a decomposition $R(G)=R(G)^{\heartsuit}\oplus R(G)^{\spadesuit}$, we obtain a decomposition
\[Z(G)=Z(G)^{\heartsuit}\oplus Z(G)^\spadesuit\] 
where $Z(G)^{\heartsuit}=\{z\in Z(G)|z|_{R(G)^{\spadesuit}}=0 \}$
and $Z(G)^{\spadesuit}=\{z\in Z(G)|z|_{R(G)^{\heartsuit}}=0 \}$.

\section{Admissibility}
We recall the notion of an admissible invariant distribution introduced by Harish-Chandra \cite[Definition 15.1]{HC99}:

\begin{definition}

\begin{enumerate}
    \item 
     An invariant distribution $f\in D(G)^G$ 
     is called 
    \emph{admissible} if for any point $x\in G$, there exists a compact open pro-$p$ subgroup 
    $U_x$ such that for any compact open $U_x'\subset U_x$
 and $(\rho,V_\rho)\in \Irr(U_x')$, we have 
\[\left(f*\Theta_\rho\right)|_{xU_x}\neq0\Rightarrow \exists g\in G\text{ such that } (V_\rho)^{U_x'\cap gU_xg^{-1}}\neq 0.\]
    Here $\Theta_\rho\in H(G)$ is the distribution character of $\rho$, normalized as a measure so that $\pi(\Theta_{\rho})$ is the projection onto $\check\rho$-isotypic components, where $\check\rho$ is the dual representation.
\item  An invariant distribution $f\in D(G)^G$ 
     is called 
$U$-\emph{admissible} with respect to a compact open pro-$p$ subgroup $U\subset G$ if  
for any compact open $U'\subset U$
 and $(\rho,V_\rho)\in \Irr(U')$, we have 
\[f*\Theta_\rho\neq0\Rightarrow \exists g\in G\text{ such that } (V_\rho)^{U'\cap gUg^{-1}}\neq 0.\]

\end{enumerate}
\end{definition}


    \begin{lemma}\label{projector}
Let $z\in Z(G)$, and let $(\rho,V_\rho)\in\Irr(K)$ be an irreducible representation of a compact open pro-$p$ subgroup $K\subset G$. Then
\[
f_z*\Theta_\rho\neq 0
\]
implies that there exists $\pi\in R(G)$ such that
\[
z(\pi)\neq 0
\qquad\text{and}\qquad
\operatorname{Hom}_K(\check\rho,\pi)\neq 0.
\]
%Here $\check\rho$ is the dual of $\rho$.
Moreover, if
\[
R(G)=\mathcal R(G)^{\heartsuit}\oplus
 R(G)^{\spadesuit}
\]
is a decomposition such that $z\in Z(G)^{\heartsuit}$, then $\pi$ may be chosen in $ R(G)^{\heartsuit}$.
\end{lemma}
\begin{proof}
Note that  $f_z*\Theta_\rho\neq0$
if and only if 
$\pi(f_z*\Theta_\rho)\neq0
\text{\ for some\ } \pi\in R(G)$.
Since \[\pi(f_z*\Theta_\rho)=z(\pi)\circ\pi(\Theta_\rho),\] we see that 
$f_z*\Theta_\rho\neq0$
implies  
 $z(\pi)\neq 0\text{ and } \pi(\Theta_\rho)\neq 0
\text{ for some } \pi\in R(G)$,
%On the other hand, \[\frac{\dim(\rho)}{\on{vol}(K)}\pi(\Theta_\rho):V_\pi\to V_\pi\] is the projection onto the $\check\rho$-isotypic component of $V_\pi$. Thus we conclude that $f_z*\Theta_\rho\neq0$ implies 
 i.e. $z(\pi)\neq 0\text{ and } \operatorname{Hom}_K(\check\rho,\pi)\neq 0 \text{\ for some\ } \pi\in R(G)$, which proves the first assertion.

Now assume $z\in Z(G)^\heartsuit$.
Then 
 we have 
\[\pi(f_z*\Theta_\rho)=\pi^\heartsuit(f_z*\Theta_\rho)\oplus\pi^\spadesuit(f_z*\Theta_\rho)\] 
where $\pi=\pi^\heartsuit\oplus\pi^\spadesuit\in R(G)=R(G)^\heartsuit\oplus R(G)^\spadesuit$.
Since  $z(\pi^\spadesuit)=0$, we have  $\pi^\spadesuit(f_z*\Theta_\rho)=z(\pi^\spadesuit)\circ\pi^\spadesuit(\Theta_\rho)=0$ and it follows that 
\[\pi^{\heartsuit}(f_z*\Theta_\rho)=\pi(f_z*\Theta_\rho)\]
and we can replace $\pi$ by the summand $\pi^\heartsuit\in R(G)^\heartsuit$
in the argument above. 
\end{proof}





Here is our key observation:

\begin{thm}\label{main}
Suppose that $U\subset G$ is a compact open pro-$p$ subgroup and $R(G)=R(G)^{\heartsuit}\oplus R(G)^{\spadesuit}$ is a
 decomposition
%with the property that there exists  a compact open pro-$p$ subgroup $U\subset G$ such that 
such that any $(\pi,V_\pi)\in R(G)^\heartsuit$ is generated by its $U$-fixed vectors $(V_\pi)^U$. Then for any $z\in Z(G)^\heartsuit$,  the corresponding 
invariant distribution $f_z$ is $U$-admissible.
 
\end{thm}
\begin{proof}
Let $\rho\in\Irr(U')$ be an irreducible representation of a compact open subgroup $U'\subset U$
such that  $f_z*\Theta_\rho\neq0$.
By Lemma \ref{projector}, there exists 
 $(\pi,V_\pi)\in R(G)^\heartsuit$ 
 satisfying $z(\pi)\neq0$ and 
 $\Hom_{U'}(\check\rho,\pi)\neq 0$.
Choose a non-zero $U'$-equivariant embedding 
$f:V_{\check\rho}\hookrightarrow V_\pi$ and let 
$0\neq w$ lie in its image. 
Since $V_\pi$ is generated by its $U$-fixed vectors, we can 
write
$w=\sum_{j=1}^n g_jv_j$ for some
$v_j\in (V_\pi)^{U}$ and $g_j\in G$. Note that we have 
$g_jv_j\in V_{\pi}^{U'\cap g_jUg_j^{-1}}$ for 
any compact open $U'\subset U$ and 
$j=1,...,n$. On the other hand, we have
\[w=\pi(\Theta_\rho)(w)=\sum_{j=1}^n\pi(\Theta_\rho)(g_jv_j)\neq 0\]
and it follows that 
\[\pi(\Theta_\rho)(g_iv_i)\neq 0\]
for some $i\in [1,n]$.
Since  $\pi(\Theta_\rho):V_\pi\to \Hom_{U'}(\check\rho,\pi)\otimes V_{\check\rho}\subset V_\pi$
is $U'$-equivariant and 
$g_iv_i$ is fixed by $U'\cap g_iUg_i^{-1}$ we conclude that 
\[0\neq\pi(\Theta_\rho)(g_iv_i)\in\Hom_{U'}(\check\rho,\pi)\otimes (V_{\check\rho})^{U'\cap g_i Ug^{-1}}.\]
Thus $(V_{\check\rho})^{U'\cap g_i Ug^{-1}}\neq0$
and, since $U'\cap g_i Ug^{-1}$ is compact
 pro-$p$ and $\ell\neq p$, it implies that 
$(V_\rho)^{U'\cap g_i Ug^{-1}}\neq0$.
\end{proof}









\section{Reductive groups}\label{reductive}
In this section we assume $G$ is a connected reductive 
group over a $p$-adic field $F$ where $\operatorname{char}(F)=0$.
%We  assume $\ell=0$ or is sufficiently large compared to $p$.

\subsection{Local integrability }
Recall the following important results of Harish-Chandra \cite{HC99}, and their generalization to mod $\ell$ coefficients in \cite{T26}:

\begin{thm}\label{HC99}
Let $f$ be an admissible invariant distribution on $G$. Then $f$ satisfies the following
\begin{enumerate}
    \item $f$ admits local character expansions as in \cite[Theorem 16.2]{HC99}.
    \item the distribution $f|_{\Grs}$
    is represented by a locally constant function 
    $F:\Grs\to k$ on the regular semisimple locus $\Grs\subset G$.
    \item When $k=\mathbb{C}$, $F$ is locally-$L^1$ on $G$ 
    and $f$ is represented by $F$. 
    \item
    When $\ell=0$ or $\ell$ is sufficiently large in the sense of \cite[Remark 3.2]{T26}, $f$ can be represented by $F$ in the sense that it can be constructed via iterated geometric series {\it loc. cit.}.
\end{enumerate}
\end{thm}
\begin{proof}
    When $k=\mathbb C$, part (1), (2) and (3) 
    are  \cite[Theorem 16.1, 16.2, 16.3]{HC99}.
    In general, as observed in \cite[Th\'{e}ore\`{m}e 5.4.4]{VW01},  the same argument in \cite{HC99} implies part (1). Part (2) is always a special case of (1) when we consider local character expansions near regular semisimple elements. Part (4) follows from Part (1) and the main result of \cite{T26}.
\end{proof}

\begin{corollary}\label{local int for z} Let $Z(G)^{\heartsuit}$ be as in Theorem \ref{main}. For any
$z\in Z(G)^\heartsuit$, the corresponding 
invariant distribution $f_z$ satisfies 
(1)-(4) of Theorem \ref{HC99}.
\end{corollary}
\begin{proof}
    Since $U$-admissibility implies admissibility, the claim follows from Theorem \ref{main} and Theorem \ref{HC99}.
\end{proof}

\subsection{Depths} For the rest of this section we assume $G$ splits over a tamely ramified extension. According to \cite{MP94,MP96,BD84} if $\ell=0$ and  \cite{V96} and \cite[Appendice A]{D09} in general, %if $\ell>0$,
for any non-negative rational number $r\in\mathbb Q_{\geq0}$, we have the depth decomposition $R(G)=R(G)_{\leq r}\oplus R(G)_{>r}$
where $R(G)_{\leq r}$
(resp. $R(G)_{>r}$)
consists of smooth representations with
irreducible subquotients of depths $\leq r$ (resp. $>r$). Let 
$K_r:=\cap_{x\in\overline{C}}G_{x,r+}$ where $\overline{C}$ is the closure of a fixed alcove.


\begin{lemma}\label{K_r} 
 The depth decomposition $R(G)=R(G)^{\heartsuit}\oplus R(G)^{\spadesuit}$, with 
   $R(G)^{\heartsuit}=R(G)_{\leq r}$
   and $R(G)^{\spadesuit}= R(G)_{>r}$, %is nice in the sense of Definition \ref{nice}, that is, there exists 
   is such that any $\pi\in R(G)^{\heartsuit}$ is generated by $K_r$-fixed vectors.
   %a compact open pro-$p$ subgroup $U\subset G$ such that    any $(\pi,V_\pi)\in R(G)_{\leq r}$ is generated by $U$-fixed vectors.
\end{lemma}
\begin{proof}
It follows from 
\cite{MP94} if $\ell=0$ and 
\cite[Lemme A.3]{D09} if $\ell>0$ that  
any $(\pi,V_\pi)\in R(G)_{\leq r}$ is generated 
by $K_r$-fixed vectors; indeed, the progenerators $P(r)$ in \cite[p. 330]{D09} are generated by $K_r$-fixed vectors by their very construction. 
%Since the family of compact open pro-$p$ subgroups of $G$ is cofinal in the set of open compact subgroups we can take $U\subset K_r$ to be any such pro-$p$ subgroup.  
\end{proof}


Consider the depth decomposition $Z(G)=Z(G)_{\leq r}\oplus Z(G)_{>r}$ of the Bernstein center
where 
$Z(G)_{\leq r}=\{z\in Z(G)|z|_{R(G)_{>r}}=0 \}$
and $Z(G)_{>r}=\{z\in Z(G)|z|_{R(G)_{\leq r}}=0 \}$.



\begin{corollary}\label{bounded depth}
      For any $z\in Z(G)_{\leq r}$
the corresponding invariant distribution 
$f_z$ satisfies 
(1)-(4) of Theorem \ref{HC99}.
\end{corollary}
\begin{proof}
It follows from Corollary \ref{local int for z} and Lemma \ref{K_r}.
\end{proof}

\subsection{Bernstein supports}
According to \cite{BD84} if $\ell=0$ and \cite[Theorem 4.22]{DHKM24} if $\ell$ is banal, there is a Bernstein decomposition $R(G)=\prod_{[M,\sigma]} R(G)_{[M,\sigma]}$ into the so-called Bernstein blocks $R(G)_{[M,\sigma]}$, where the index set $[M,\sigma]$ is the inertial equivalence classes of supercuspidal supports. We say that $z\in Z(G)$ has bounded Bernstein support if its restriction to all but finitely many Bernstein blocks is zero.

\begin{corollary}\label{supp}\label{bounded support} Suppose $\ell=0$ or $\ell$ is banal. For any $z\in Z(G)$ with bounded Bernstein support, 
the corresponding invariant distribution 
$f_z$ satisfies 
(1)-(4) of Theorem \ref{HC99}.
\end{corollary}
\begin{proof}
Since unramified twists  and parabolic induction preserve depth,
we have 
$R(G)_{[M,\sigma]}\subset R(G)_{\leq r}$ for some $r$.
Thus any $z\in Z(G)$ with bounded Bernstein support lies in $Z(G)_{\leq r}$ for some sufficiently large $r$ and the claim follows from Corollary \ref{bounded depth}.
\end{proof}
\begin{remark}
    In the case $\ell=0$, Corollary \ref{bounded depth} or \ref{supp} above provides an alternative argument for the main results in \cite{MT02}. 
    The case $\ell>0$ appears to be new and does not follow directly from the methods of \emph{loc. cit.}
\end{remark}

\begin{remark}
Consider the case where $G$ is a reductive group over a local function field $F$ of characteristic $p>0$. We expect Harish-Chandra's Theorem (Theorem \ref{HC99}(1)) to remain valid when $p$ is sufficiently large. When this is the case, all results in Section \ref{reductive} follow. 
Moreover, with some modifications, we expect that the results of the note remain valid when $k$ is replaced by any $\mathbb Z[1/p]$-algebra.


\end{remark}

\subsection*{Acknowledgments} We thank the NCTS-National Center for Theoretical Sciences at Taipei where parts of this work were done. The research of T.-H. Chen is supported by NSF grant DMS-214372
and Simons Fellowships.
The research of C.-C. Tsai is supported by NSTC grant 115-2628-M-001-002.


\begin{thebibliography}{9999}
%\bibitem[AGKRV]{AGKRV} D. Gaitsgory, D. Kazhdan, N. Rozenblyum, Y. Varshavsky,
%A toy model for the Drinfeld-Lafforgue shtuka construction, arXiv:1908.05420.



\bibitem[BD84]{BD84}
J. Bernstein and P. Deligne, “Le ‘centre’ de Bernstein,” in Representations des groupes réductifs sur un corps local, Hermann, Paris, 1984, pp. 1–32




\bibitem[D09]{D09} J.-F. Dat, Finitude pour les repr\'esentations lisses de groupes $p$-adiques, J. Inst. Math. Jussieu {\bf 8} (2009), no.~2, 261--333; MR2485794

\bibitem[DHKM24]{DHKM24} 
J.-F. Dat, D. Helm, R. Kurinczuk, and G. Moss, Local Langlands in families: The banal case. Preprint,
arXiv:2406.09283 [math.RT] (2024), 2024, https://arxiv.org/abs/2406.09283.

\bibitem[HC99]{HC99} 
Harish-Chandra,
Admissible invariant distributions on reductive $p$-adic groups, University Lecture Series, vol. 16, American Mathematical Society, Providence, RI, 1999, Preface and notes by Stephen DeBacker and Paul J. Sally, Jr..
%Lie Theories and Their Applications, Queen's Papers in Pure and Applied Mathematics 48 (1978), 281–347.

\bibitem[MP94]{MP94} 
A. Moy and G. Prasad, Unrefined minimal K-types for p-adic groups, Invent. Math.
116 (1994), no. 1-3, 393–408. MR 1253198 (95f:22023)

\bibitem[MP96]{MP96} 
A. Moy, G. Prasad,  Jacquet functors and unrefined minimal K-types. Commentarii Mathematici Helvetici 71, 98–121 (1996). https://doi.org/10.1007/BF02566411

\bibitem[MT02]{MT02} 
A. Moy and M. Tadi\'{c},
 The Bernstein center in terms of invariant locally integrable functions,
Representation Theory 6 (2002), 313–329.
DOI: 10.1090/S1088-4165-02-00181-4





\bibitem[T26]{T26} C.-C. Tsai, On local integrability results for $p$-adic reductive groups, \url{https://arxiv.org/abs/2604.15079v1},


\bibitem[V96]{V96} 
M.F. Vign\'{e}ras, Representations $\ell$-modulaires d’un groupe reductif p-adique avec
$\ell\neq p$, Progress in Mathematics, vol. 137, Birkh¨auser Boston, Inc., Boston, MA, 1996.
MR 1395151

\bibitem[VW01]{VW01} 
M.-F. Vign\'{e}ras and J.-L. Waldspurger, “Premiers réguliers de l’analyse harmonique mod $\ell$ d’un groupe réductif $p$-adique,” Journal für die reine und angewandte Mathematik 535 (2001), 165–205



\end{thebibliography}
\end{document}
